\chapter{Installation and the First Workflow}
\label{ch:getting-started}

This chapter takes a new user from an empty machine to a verified, running
DAGonStar workflow, following the project's own \texttt{getting\_started.md}
guide.

\section{Requirements}

\begin{itemize}[leftmargin=1.4em]
  \item Python 3.8 or newer.
  \item A POSIX-like shell for the included command examples.
  \item Git.
\end{itemize}

Optional integrations require additional, separately verified services:

\begin{itemize}[leftmargin=1.4em]
  \item Docker, for Docker tasks (\Cref{ch:backends});
  \item SSH access, for remote tasks;
  \item Slurm, for HPC scheduler tasks;
  \item cloud provider credentials, for cloud tasks;
  \item Globus or SKYCDS credentials, for specialized data movement.
\end{itemize}

\section{Installing DAGonStar from Source}

\begin{lstlisting}[style=bash, caption={Full development installation.}]
git clone https://github.com/DagOnStar/dagonstar.git
cd dagonstar
python3 -m venv .venv
. .venv/bin/activate
python -m pip install --upgrade pip
python -m pip install -r requirements.txt
python -m pip install -e .
\end{lstlisting}

For a smaller installation footprint, install only the core package and the
optional integrations actually needed:

\begin{lstlisting}[style=bash, caption={Minimal install with selective extras.}]
python -m pip install -e .
python -m pip install -e ".[docker]"   # Docker tasks
python -m pip install -e ".[cloud]"    # cloud tasks
python -m pip install -e ".[globus]"   # Globus staging
python -m pip install -e ".[api]"      # Flask workflow service
\end{lstlisting}

\begin{practicebox}
This extras-based packaging is itself a reproducibility decision: a published
workflow's installation instructions should record \emph{exactly} which
extras were required, so a future reader does not need to guess whether
Docker or cloud credentials are necessary to reproduce a given experiment.
\end{practicebox}

\section{Creating a Configuration File}

\begin{lstlisting}[style=bash]
cp dagon.ini.sample dagon.ini
\end{lstlisting}

For a first local workflow, the sample defaults suffice. Confirm the two keys
that matter most for local execution:

\begin{lstlisting}[style=ini]
[dagon_service]
use = False

[batch]
scratch_dir_base=/tmp/
\end{lstlisting}

The full configuration model is the subject of \Cref{ch:configuration} and
Appendix~C.

\section{Writing a First Workflow}

\begin{lstlisting}[caption={\texttt{first\_workflow.py}: a two-task, data-oriented workflow.}]
from dagon import Workflow
from dagon.task import DagonTask, TaskType


if __name__ == "__main__":
    workflow = Workflow("FirstWorkflow")

    prepare = DagonTask(
        TaskType.BATCH,
        "prepare",
        "mkdir -p data; echo 'temperature,depth' > data/profile.csv; "
        "echo '18.5,12' >> data/profile.csv",
    )

    analyze = DagonTask(
        TaskType.BATCH,
        "analyze",
        "wc -l workflow:///prepare/data/profile.csv > summary.txt",
    )

    workflow.add_task(prepare)
    workflow.add_task(analyze)
    workflow.make_dependencies()
    workflow.run()
\end{lstlisting}

\begin{lstlisting}[style=bash]
python first_workflow.py
\end{lstlisting}

\subsection{What Happens, Step by Step}

\begin{enumerate}[leftmargin=1.6em]
  \item \texttt{prepare} creates a small CSV file inside its own scratch
    directory.
  \item \texttt{analyze} references that file with
    \texttt{workflow:///prepare/data/profile.csv}.
  \item \texttt{workflow.make\_dependencies()} detects, by parsing
    \texttt{analyze}'s command, that \texttt{analyze} depends on
    \texttt{prepare}.
  \item DAGonStar stages the referenced file into the \texttt{analyze}
    task's working directory.
  \item \texttt{summary.txt} is produced by the \texttt{analyze} task.
\end{enumerate}

This five-step trace is a miniature instance of the fifteen-step lifecycle
formalized in \Cref{sec:execution-lifecycle}; the reader is encouraged to map
each numbered step here onto the corresponding lifecycle phase.

\section{Inspecting Generated Task Directories}

By default, task directories are created under \texttt{/tmp/}, with names
that include a millisecond timestamp and the task name, e.g.:

\begin{lstlisting}[style=bash]
/tmp/1710000000000-prepare
/tmp/1710000000100-analyze
\end{lstlisting}

Each task directory contains:

\begin{itemize}[leftmargin=1.4em]
  \item task output files;
  \item a \texttt{.dagon/} directory;
  \item generated launcher scripts;
  \item captured standard output.
\end{itemize}

\Cref{ch:lessons-03-04} examines the contents of \texttt{.dagon/} in detail,
treating the generated launcher as a first-class provenance artifact.

\section{Running the Unit Test Suite}

\begin{lstlisting}[style=bash]
python -m unittest discover -s tests -v
\end{lstlisting}

The test suite verifies behavior that the project commits to keeping stable
while new features are developed; running it after installation is a useful
sanity check that the local environment matches the documented one.

\section{Understanding the Installation Footprint}

The two installation paths given above --- the full
\texttt{requirements.txt} install and the minimal extras-based install ---
are not merely a convenience choice; they encode a deliberate packaging
philosophy that is worth making explicit before proceeding further, since it
recurs throughout Part~II whenever a new execution backend is introduced.

The core package --- \texttt{Workflow}, \texttt{DagonTask}, the
\texttt{BATCH} backend, dependency resolution, and local data staging ---
has, by design, no dependency on Docker, any cloud SDK, Globus, or a web
framework. A researcher whose entire pipeline consists of local shell
commands and \texttt{workflow://} references can install DAGonStar with a
minimal footprint and a correspondingly small attack surface and dependency
tree. Every additional capability examined in \Cref{ch:backends} onward ---
container execution, cloud provisioning, Globus-mediated transfer, the
optional HTTP service --- is gated behind an explicit \texttt{pip} extra,
and, as \Cref{sec:known-constraints} noted, the corresponding Python modules
use \emph{lazy imports}, meaning the module is not even imported, let alone
required to be installed, unless a workflow actually constructs a task of
that type.

\begin{practicebox}
This has a direct, practical consequence for reproducibility reporting,
foreshadowed already in this chapter's earlier discussion of extras: a
workflow's installation instructions should state not just ``install
DAGonStar,'' but the precise set of extras required, because that set is
itself a statement about which execution backends --- and therefore which
external services and credentials --- a reader will need to provision before
they can run the workflow end to end. A workflow using only \texttt{BATCH}
and \texttt{LLM} tasks against a mock provider, for instance, needs no
extras at all; one that also schedules \texttt{DOCKER} tasks needs the
\texttt{docker} extra and a running daemon.
\end{practicebox}

\section{The First Workflow as a Minimal Reproducible Experiment}

It is worth pausing on \texttt{first\_workflow.py} a second time, after the
step-by-step trace already given, to read it specifically as an instance of
what \Cref{ch:introduction} called a scientific benefit rather than merely
as a syntactic example.

Consider what would be required to reproduce this two-line ``experiment''
--- the claim that a small CSV profile, once line-counted, yields a
particular summary --- using only shell history. A second researcher would
need to know, without being told explicitly, that the \texttt{wc -l}
invocation depended on the CSV having already been generated by an earlier,
separately run command; nothing in the shell history itself records that
dependency as a structural fact, only as an accident of the order in which
commands happened to be typed. The DAGonStar version of the same experiment
makes the dependency a first-class, parseable part of the program text:
\texttt{workflow:///prepare/data/profile.csv} cannot be misread or silently
decoupled from its producer, because \texttt{make\_dependencies()} re-derives
the edge from the reference every time the script runs, on any machine, in
any order the engine chooses to schedule it. This is precisely the
\emph{traceability} benefit named in \Cref{ch:introduction}, made concrete
in the smallest possible example before Part~I's later chapters scale it up
to checkpointed, staged, multi-backend pipelines.

\begin{exercise}
Modify \texttt{first\_workflow.py} so that a third task,
\texttt{validate}, depends on \texttt{analyze}'s \texttt{summary.txt} output
and exits with a non-zero status if the line count it contains is zero
(indicating an empty or malformed profile). Run the modified workflow and
verify, using the techniques of \Cref{sec:execution-lifecycle}, that the
\texttt{validate} task is correctly recorded as \texttt{FAILED} when you
intentionally corrupt the \texttt{prepare} task's output to produce an empty
file.
\end{exercise}

\section{Common Early Pitfalls}

New users of any workflow engine tend to encounter a small, predictable set
of early mistakes; naming them explicitly here, before Part~I develops the
underlying mechanisms in full, should save the reader some early debugging
time.

\begin{description}[leftmargin=1.8em, style=nextline]
  \item[Forgetting \texttt{make\_dependencies()} while still inspecting the
    graph.] As \Cref{ch:meta-async} will show, \texttt{run()} and
    \texttt{launch()} both build missing dependencies automatically, so
    omitting the explicit call does not break execution --- but it does mean
    that any code placed \emph{between} task construction and
    \texttt{run()} that inspects \texttt{task.prevs} or \texttt{task.nexts}
    will see an empty, not-yet-resolved graph. Call
    \texttt{make\_dependencies()} explicitly whenever the graph must be
    inspected or validated before execution.
  \item[Confusing the producer task's working directory with the consumer's.]
    A \texttt{workflow://} path is always relative to the \emph{producer}
    task's directory (\Cref{ch:schema}), never to the consumer's; a common
    early error is writing a reference as though the path were relative to
    wherever the consumer happens to be running.
  \item[Assuming sibling tasks run in a particular order.] As
    \Cref{sec:topological-ordering} makes precise, two tasks with no
    dependency edge between them have no guaranteed relative order; code
    that implicitly relies on one finishing before the other without an
    edge connecting them is a latent bug that may not surface until the
    workflow is run on a machine with different thread-scheduling
    behavior.
  \item[Treating \texttt{/tmp} as durable storage.] The default
    \texttt{scratch\_dir\_base} of \texttt{/tmp/} is convenient for the
    tutorial-scale examples in this book, but \texttt{/tmp} is frequently
    cleared on system reboot or by periodic cleanup policies; any workflow
    relying on checkpoint reuse (\Cref{ch:lessons-05-07}) across a span of
    time longer than a single working session should configure a more
    durable \texttt{scratch\_dir\_base}, as \Cref{ch:configuration}
    discusses.
\end{description}

\section{Where to Go Next}

\begin{itemize}[leftmargin=1.4em]
  \item \Cref{ch:configuration} develops the complete configuration model.
  \item \Cref{ch:schema} formalizes the \texttt{workflow://} schema used
    implicitly in this chapter's example.
  \item \Cref{ch:lessons-01-02} onward follow the course's 48 lectures, each
    with a tested example in \texttt{examples/}.
  \item \Cref{ch:external-software} explains how to wrap compiled,
    command-line scientific models such as WRF.
  \item \Cref{ch:architecture} should be revisited before modifying engine
    internals.
\end{itemize}
